\section*{Capítulo \ref{ch:g7:stats} --- Organizando dados}

\begin{solution}{exo:g7:stats:1}
Contando cada nota:
\begin{center}
\renewcommand{\arraystretch}{1.2}
\begin{tabular}{l|ccccc|c}
nota & $9$ & $11$ & $12$ & $14$ & $15$ & total \\
\hline
freq. absoluta & $4$ & $3$ & $7$ & $3$ & $3$ & $20$ \\
\end{tabular}
\end{center}
\end{solution}

\begin{solution}{exo:g7:stats:2}
\omterm{def:g7:stats:frequency}{Frequências relativas}: $\frac{4}{20} = 20\,\%$,
$\frac{3}{20} = 15\,\%$, $\frac{7}{20} = 35\,\%$, $15\,\%$, $15\,\%$.
\omterm{def:g1:addition:def}{Soma}: $20 + 15 + 35 + 15 + 15 = 100\,\%$. \checkmark
\end{solution}

\begin{solution}{exo:g7:stats:3}
Primeira pesquisa: $\frac{30}{50} = 0.6 = 60\,\%$. Segunda:
$\frac{45}{90} = 0.5 = 50\,\%$. A primeira pesquisa mostra a proporção
maior, embora tenha menos famílias com um carro em número absoluto.
\end{solution}

\begin{solution}{exo:g7:stats:4}
$\intco{5}{15}$: $8, 12, 9, 14, 11$, frequência $5$.

$\intco{15}{25}$: $23, 22, 18, 16$, frequência $4$.

$\intco{25}{35}$: $31, 27, 29, 25, 33$, frequência $5$.

$\intco{35}{45}$: $35$, frequência $1$. Total $15$. \checkmark
\end{solution}

\begin{solution}{exo:g7:stats:5}
Quatro barras de alturas $8$, $12$, $6$, $4$ sobre os valores $0$, $1$,
$2$, $3$ (graduação de $2$ em $2$ funciona bem).
\end{solution}

\begin{solution}{exo:g7:stats:6}
Total: $30$ famílias. \omterm{def:g7:stats:frequency}{Frequências relativas}: $\frac{8}{30}$,
$\frac{12}{30}$, $\frac{6}{30}$, $\frac{4}{30}$. Ângulos
($\times 360^\circ$): $96^\circ$, $144^\circ$, $72^\circ$, $48^\circ$;
\omterm{def:g1:addition:def}{soma} $96 + 144 + 72 + 48 = 360^\circ$. \checkmark
\end{solution}

\begin{solution}{exo:g7:stats:7}
Porcentagens: verão $\frac{162}{360} = 45\,\%$; primavera
$\frac{90}{360} = 25\,\%$; outono e inverno $\frac{54}{360} = 15\,\%$
cada. Entre $60$ pessoas, o verão foi escolhido por
$0.45 \times 60 = 27$ delas.
\end{solution}

\begin{solution}{exo:g7:stats:8}
Turma A: $\frac{12}{30} = 40\,\%$. Turma B: $\frac{14}{40} = 35\,\%$. A
turma A tem a maior proporção de alunos que vão a pé, mesmo a turma B
tendo mais alunos a pé em número absoluto.
\end{solution}

\begin{solution}{exo:g7:stats:9}
Com o eixo começando em $9\,700$, as três barras têm alturas visíveis
de $100$, $300$ e $400$ unidades: a última parece quatro vezes a
primeira, sugerindo que as vendas quadruplicaram. Com um eixo a partir
de $0$, as barras ficam quase iguais (de $9\,800$ a $10\,100$ é uma
alta de cerca de $3\,\%$): a imagem honesta mostra vendas quase
estáveis.
\end{solution}

\begin{solution}{exo:g7:stats:10}
As \omterm{def:g7:stats:frequency}{frequências relativas} somam $1$: $f = 1 - 0.35 - 0.4 = 0.25$.
\omterm{def:g7:stats:frequency}{Frequências absolutas} (total $80$): $0.35 \times 80 = 28$;
$0.4 \times 80 = 32$; $0.25 \times 80 = 20$. Confere:
$28 + 32 + 20 = 80$.
\end{solution}

\begin{solution}{exo:g7:stats:11}
Meninas: $55\,\%$ de $400 = 220$; meninos: $180$. Meninas que almoçam
no refeitório: $0.4 \times 220 = 88$; meninos: $0.6 \times 180 = 108$.
Total: $88 + 108 = 196$ alunos, isto é, $\frac{196}{400} = 49\,\%$ da
escola --- entre $40\,\%$ e $60\,\%$, mais perto da taxa das meninas
porque elas são mais numerosas.
\end{solution}

\begin{solution}{pb:g7:stats:1}
\textbf{1.} \omterm{def:g7:stats:frequency}{Frequências absolutas}: $120 > 90$, a escola A ganha.
\omterm{def:g7:stats:frequency}{Frequências relativas}: $\frac{120}{400} = 0.30 = 30\,\%$ contra
$\frac{90}{250} = 0.36 = 36\,\%$: a escola B ganha. O prêmio deve
seguir a \omterm{def:g7:stats:frequency}{frequência relativa} --- a escola B recicla uma parcela maior
do que usa; a escola A apenas bebe mais suco.

\textbf{2.} Professores: $\frac{18}{30} = 0.60 = 60\,\%$. Alunos:
$\frac{210}{600} = 0.35 = 35\,\%$. Juntos:
$\frac{18 + 210}{630} = \frac{228}{630} \approx 0.36 = 36\,\%$. Os
alunos são vinte vezes mais numerosos, então a frequência combinada é
puxada quase inteiramente para o lado deles --- o grupo grande carrega
o peso grande.

\textbf{3.} As \omterm{def:g7:stats:frequency}{frequências relativas} têm de somar $100\,\%$
(\cref{def:g7:stats:frequency}), mas $45 + 30 + 20 + 10 = 105$: pelo
menos um setor está errado.

\textbf{4.} \omterm{def:g7:stats:frequency}{Frequência absoluta} correspondente a $0.15$:
$0.15 \times 40 = 6$. As frequências até aqui somam
$14 + 10 + 6 = 30$, então a última é $40 - 30 = 10$, com \omterm{def:g7:stats:frequency}{frequência
relativa} $\frac{10}{40} = 0.25$. (Confere:
$0.35 + 0.25 + 0.15 + 0.25 = 1$.)

\textbf{5.} Ângulos: $0.40 \times 360 = 144^\circ$,
$0.35 \times 360 = 126^\circ$, $0.25 \times 360 = 90^\circ$; e
$144 + 126 + 90 = 360^\circ$: a pizza fecha.

\textbf{6.} $\frac{2.4}{10} = 0.24 = 24\,\%$ das cédulas voltaram.

\textbf{7.} As listas telefônicas e os registros de automóveis
relacionavam sobretudo lares ricos, cujo voto era diferente do voto do
país; os $2.4$ milhões de respostas foram uma contagem gigantesca
tirada da população errada. Como na questão 1: o que importa não é
quantos você conta, e sim quem --- uma frequência calculada numa
amostra distorcida descreve a amostra, não o país.

\textbf{8.} Apoio verdadeiro:
$0.7 \times 1\,000 + 0.3 \times 9\,000 = 700 + 2\,700 = 3\,400$
apoiadores em $10\,000$: $34\,\%$. A pesquisa:
$0.7 \times 500 + 0.3 \times 500 = 350 + 150 = 500$ em $1\,000$:
$50\,\%$ --- dezesseis pontos acima, porque os ricos são \omterm{def:g3:division:half}{metade} da
amostra, mas um décimo da cidade.

\textbf{9.} Respostas: $0.4 \times 200 = 80$ insatisfeitos e
$0.1 \times 800 = 80$ satisfeitos, então $160$ respostas, das quais
$80$ insatisfeitas: insatisfação medida $\frac{80}{160} = 50\,\%$.
Insatisfação verdadeira: $\frac{200}{1\,000} = 20\,\%$. Ninguém mentiu
--- os insatisfeitos simplesmente respondem mais, e a pesquisa os ouve
mais alto.

\textbf{10.} Proporções verdadeiras: um décimo de ricos, ou seja, $100$
entrevistas com ricos e $900$ com modestos:
$0.7 \times 100 + 0.3 \times 900 = 70 + 270 = 340$ em $1\,000$: a
pesquisa consertada prevê $34\,\%$ --- a verdade da questão 8.
Cinquenta mil entrevistas bem escolhidas vencem dois milhões mal
escolhidas.

\textbf{11.} Hospital A: leves $\frac{90}{100} = 90\,\%$; graves
$\frac{280}{400} = 70\,\%$; no total
$\frac{90 + 280}{500} = \frac{370}{500} = 74\,\%$.

\textbf{12.} Hospital B: leves $\frac{340}{400} = 85\,\%$; graves
$\frac{65}{100} = 65\,\%$; no total
$\frac{340 + 65}{500} = \frac{405}{500} = 81\,\%$. O hospital A vence
nos casos leves ($90 > 85$) \emph{e} nos casos graves ($70 > 65$).

\textbf{13.} No total, B mostra $81\,\%$ contra os $74\,\%$ de A: o
hospital que perde em \emph{todas} as categorias ganha no geral. (Essa
inversão tem nome: paradoxo de Simpson.)

\textbf{14.} As misturas são opostas: os pacientes de A são na maioria
graves ($400$ de $500$), e os de B, na maioria leves ($400$ de $500$).
Casos graves se curam menos onde quer que sejam tratados, então o
número geral de A é puxado para baixo pelos casos difíceis que ele
aceita --- o total é uma mistura ponderada, e os pesos diferem entre os
hospitais. Um paciente grave deve escolher o hospital A: $70\,\%$ vence
$65\,\%$ na única linha que lhe diz respeito.

\textbf{15.} Por exemplo: ``A sua comparação mistura pacientes leves e
graves, e os dois hospitais tratam misturas opostas --- separados por
gravidade, o hospital A cura uma parcela maior dos \emph{dois} tipos.
Números gerais só podem ser comparados quando os grupos por trás deles
são parecidos; caso contrário, quem decide o vencedor são os pesos, e
não a qualidade.''
\end{solution}
